\documentclass[10pt,a4paper]{amsart}
\usepackage{amsmath,amssymb,amsthm,graphicx}
\usepackage{verbatim}
\newenvironment{reference}{\comment}{\endcomment}
%\newenvironment{reference}{}{}
\newenvironment{slogan}{\comment}{\endcomment}
\newenvironment{history}{\comment}{\endcomment}

% For commutative diagrams we use Xy-pic
\usepackage[all]{xy}

% We use 2cell for 2-commutative diagrams.
\xyoption{2cell}
% Theorem environments.
%
\theoremstyle{plain}
\newtheorem{theorem}[subsection]{Theorem}
\newtheorem{proposition}[subsection]{Proposition}
\newtheorem{lemma}[subsection]{Lemma}

\theoremstyle{definition}
\newtheorem{definition}[subsection]{Definition}
\newtheorem{example}[subsection]{Example}
\newtheorem{exercise}[subsection]{Exercise}
\newtheorem{situation}[subsection]{Situation}

\theoremstyle{remark}
\newtheorem{remark}[subsection]{Remark}
\newtheorem{remarks}[subsection]{Remarks}

\numberwithin{equation}{subsection}

% Macros
%
\def\lim{\mathop{\mathrm{lim}}\nolimits}
\def\colim{\mathop{\mathrm{colim}}\nolimits}
\def\Spec{\mathop{\mathrm{Spec}}}
\def\Hom{\mathop{\mathrm{Hom}}\nolimits}
\def\Ext{\mathop{\mathrm{Ext}}\nolimits}
\def\SheafHom{\mathop{\mathcal{H}\!\mathit{om}}\nolimits}
\def\SheafExt{\mathop{\mathcal{E}\!\mathit{xt}}\nolimits}
\def\Sch{\mathit{Sch}}
\def\Mor{\mathop{\mathrm{Mor}}\nolimits}
\def\Ob{\mathop{\mathrm{Ob}}\nolimits}
\def\Sh{\mathop{\mathit{Sh}}\nolimits}
\def\NL{\mathop{N\!L}\nolimits}
\def\CH{\mathop{\mathrm{CH}}\nolimits}
\def\proetale{{pro\text{-}\acute{e}tale}}
\def\etale{{\acute{e}tale}}
\def\QCoh{\mathit{QCoh}}
\def\Ker{\mathop{\mathrm{Ker}}}
\def\Im{\mathop{\mathrm{Im}}}
\def\Coker{\mathop{\mathrm{Coker}}}
\def\Coim{\mathop{\mathrm{Coim}}}

% Boxtimes
%
\DeclareMathSymbol{\boxtimes}{\mathbin}{AMSa}{"02}

%
% Macros for moduli stacks/spaces
%
\def\QCohstack{\mathcal{QC}\!\mathit{oh}}
\def\Cohstack{\mathcal{C}\!\mathit{oh}}
\def\Spacesstack{\mathcal{S}\!\mathit{paces}}
\def\Quotfunctor{\mathrm{Quot}}
\def\Hilbfunctor{\mathrm{Hilb}}
\def\Curvesstack{\mathcal{C}\!\mathit{urves}}
\def\Polarizedstack{\mathcal{P}\!\mathit{olarized}}
\def\Complexesstack{\mathcal{C}\!\mathit{omplexes}}
% \Pic is the operator that assigns to X its picard group, usage \Pic(X)
% \Picardstack_{X/B} denotes the Picard stack of X over B
% \Picardfunctor_{X/B} denotes the Picard functor of X over B
\def\Pic{\mathop{\mathrm{Pic}}\nolimits}
\def\Picardstack{\mathcal{P}\!\mathit{ic}}
\def\Picardfunctor{\mathrm{Pic}}
\def\Deformationcategory{\mathcal{D}\!\mathit{ef}}
\usepackage{cleveref}
\crefname{originallemma}{original lemma}{original lemmas}
\Crefname{originallemma}{Original lemma}{Original lemmas}
\makeatletter
\newlabel{algebra-lemma-localization-catenary}{{00NJ}{source}{}{sourcepointer-1}{}}
\newlabel{algebra-lemma-localization-catenary@cref}{{[originallemma][0][]00NJ}{[1][0][]source}{}{}{}}
\makeatother
\crefname{originallemma}{original lemma}{original lemmas}
\Crefname{originallemma}{Original lemma}{Original lemmas}
\makeatletter
\newlabel{algebra-lemma-noetherian-dim-1-Jacobson}{{00KX}{source}{}{sourcepointer-2}{}}
\newlabel{algebra-lemma-noetherian-dim-1-Jacobson@cref}{{[originallemma][0][]00KX}{[1][0][]source}{}{}{}}
\makeatother
\crefname{originallemma}{original lemma}{original lemmas}
\Crefname{originallemma}{Original lemma}{Original lemmas}
\makeatletter
\newlabel{morphisms-lemma-ubiquity-uc}{{02JB}{source}{}{sourcepointer-3}{}}
\newlabel{morphisms-lemma-ubiquity-uc@cref}{{[originallemma][0][]02JB}{[1][0][]source}{}{}{}}
\makeatother
\crefname{originallemma}{original lemma}{original lemmas}
\Crefname{originallemma}{Original lemma}{Original lemmas}
\makeatletter
\newlabel{varieties-lemma-complement-codim-1-closed-points}{{09NA}{source}{}{sourcepointer-4}{}}
\newlabel{varieties-lemma-complement-codim-1-closed-points@cref}{{[originallemma][0][]09NA}{[1][0][]source}{}{}{}}
\makeatother
\title{A Jacobson universally catenary Noetherian domain with maximal ideals of local dimension 1 and 2}
\author{The Stacks Project Authors}
\date{stacks-project:examples.tex:local-snapshot-20261012; GNU FDL 1.2 or later}
\begin{document}
\maketitle
\noindent\textit{Editorial note.} All delivered mathematical text is the author's own, copied from the pinned Stacks Project chapter file examples.tex; no statement or proof step was rewritten, reordered or removed.\par\medskip
\noindent\textit{Editorial note.} The counted proof is the author's complete argument: The counted proof delegates everything to the context construction and adds the two universal-catenary citations (Algebra, Morphisms - both intra-corpus); the author placeholder '(insert future reference here)' in the Claim is inside the context span and is printed verbatim.\par\medskip
\noindent\textit{Editorial note.} The label lemma-Noetherian-Jacobson names the counted statement, and the proof environment immediately following it in the source is delivered verbatim as its proof.\par\medskip
\noindent\textit{Editorial note.} The counted statement is carried by the author's own lemma environment in the chapter source.\par\medskip
\noindent\textit{Allowed standard prerequisites (editorial).}\begin{itemize}
\item Context span 1848-1917: the construction - k algebraic closure of a finite field, A = k[x,y], C = V(x) with enumerated closed points, auxiliary curves C\_i, the Claim producing irreducible Z\_n with \{p\_n\} = Z\_n cap (C cup C\_1 cup ... cup C\_n) via Cartier divisors and torsion Picard group of an affine curve, then S generated by the f\_n with Z\_n = V(f\_n) and B = S\textasciicircum{}\{-1\}A, with the analysis of which maximal ideals survive (infinitely many from A, plus xB) yielding local rings of dimensions 1 and 2
\item Universally catenary rings: Algebra, lemma-localization-catenary and Morphisms, lemma-ubiquity-uc
\end{itemize}\medskip
\ifdefined\hypertarget\hypertarget{sourcepointer-1}{}\fi
\noindent\textit{Source reference.} Original lemma 00NJ: sources/stacks-chow/tags.txt, source lines 861--861. Prerequisite pointer invoked at 1928; it names a standard fact of a sibling chapter that the counted argument uses, and every step of the counted proof built on it is retained verbatim in the delivered body.\par\medskip
\ifdefined\hypertarget\hypertarget{sourcepointer-2}{}\fi
\noindent\textit{Source reference.} Original lemma 00KX: sources/stacks-chow/tags.txt, source lines 767--767. Prerequisite pointer invoked at 1995; it names a standard fact of a sibling chapter that the counted argument uses, and every step of the counted proof built on it is retained verbatim in the delivered body.\par\medskip
\ifdefined\hypertarget\hypertarget{sourcepointer-3}{}\fi
\noindent\textit{Source reference.} Original lemma 02JB: sources/stacks-chow/tags.txt, source lines 3330--3330. Prerequisite pointer invoked at 1929; it names a standard fact of a sibling chapter that the counted argument uses, and every step of the counted proof built on it is retained verbatim in the delivered body.\par\medskip
\ifdefined\hypertarget\hypertarget{sourcepointer-4}{}\fi
\noindent\textit{Source reference.} Original lemma 09NA: sources/stacks-chow/tags.txt, source lines 12208--12208. Prerequisite pointer invoked at 1959; it names a standard fact of a sibling chapter that the counted argument uses, and every step of the counted proof built on it is retained verbatim in the delivered body.\par\medskip
Let $k$ be the algebraic closure of a finite field.
Let $A = k[x, y]$ and $X = \Spec(A)$.
Let $C = V(x)$ be the $y$-axis (this could be any other
$1$-dimensional integral closed subscheme of $X$).
Let $C_1, C_2, C_3, \ldots$ be an enumeration
of the other integral closed subschemes of $X$ of dimension $1$.
Let $p_1, p_2, p_3, \ldots$ be an
enumeration of the closed points of $C$.

\medskip\noindent
Claim: for every $n$ there exists an irreducible closed
$Z_n \subset X$ of dimension $1$ such that
$$
\{p_n\} = Z_n \cap (C \cup C_1 \cup \ldots \cup C_n)
$$
set theoretically. To do this set
$Y = C \cup C_1 \cup C_2 \cup \ldots \cup C_n$.
This is a reduced affine algebraic scheme of dimension $1$ over $k$.
It is enough to find $f \in k[x, y]$ with $V(f) \cap  Y = \{p_n\}$
set theoretically because then we can take $Z_n$ to be a suitable
irreducible component of $V(f)$. Since the restriction map
$$
k[x, y] \longrightarrow \Gamma(Y, \mathcal{O}_Y)
$$
is surjective, it suffices to find a regular function $g$ on $Y$
whose zero set is $\{p_n\}$ set theoretically.
To see this is possible, we choose an effective Cartier divisor
$D \subset Y$ whose support is $p_n$ (this is possible by
Varieties, Lemma \ref{varieties-lemma-complement-codim-1-closed-points}).
Thus it suffices to show that $\mathcal{O}_X(ND) \cong \mathcal{O}_X$
for some $N > 0$. But the Picard group of an affine $1$-dimensional
algebraic scheme over the algebraic closure of a finite field is torsion
(insert future reference here) and we conclude the claim is true.

\medskip\noindent
Choose $Z_n$ as above for all $n$. Since $k[x, y]$ is a UFD
we may write $Z_n = V(f_n)$ for some irreducible element $f_n \in A$.
Let $S \subset k[x, y]$ be the multiplicative subset generated
by $f_1, f_2, f_3, \ldots$. Consider the Noetherian ring $B = S^{-1}A$.

\medskip\noindent
Obviously, the ring map $A \to B$ identifies local rings
and induces an injection $\Spec(B) \to \Spec(A)$.
Moreover, looking at the curve $C_1$ we see that
only the points of $C \cap C_1$ are removed when passing
from $\Spec(A)$ to $\Spec(B)$.
In particular, we see that $\Spec(B)$ has an infinite
number of maximal ideals corresponding to maximal
ideals of $A$. On the other hand, $xB$ is a maximal ideal because
the spectrum of $B/xB$ consists of a unique prime ideal
as we removed all the closed points of $C = V(x)$ (but not
the generic point).
Finally, for $i \geq 1$ consider the curve $C_i$.
Write $C_i = V(g_i)$ for $g_i \in A$ irreducible.
If $C_i = Z_n$ for some $n$, then $g_iB$ is the unit ideal.
If not, then all but finitely many of the closed points
of $C_i$ survive the passage from $A$ to $B$:
namely, only the points of
$(Z_1 \cup \ldots \cup Z_{i - 1} \cup C) \cap C_i$
are removed from $C_i$.

\medskip\noindent
The structure of the prime spectrum of $B$ given above
shows that $B$ is Jacobson by
Algebra, Lemma \ref{algebra-lemma-noetherian-dim-1-Jacobson}.
The maximal ideals are the maximal ideals of $A$
which are in $\Spec(B)$ (and there an inifinitude of these)
together with the maximal ideal $xB$. Thus we see that we have
local rings of dimensions $1$ and $2$.


\begin{lemma}
\label{lemma-Noetherian-Jacobson}
There exists a Jacobson, universally catenary, Noetherian domain $B$
with maximal ideals $\mathfrak m_1, \mathfrak m_2$ such that
$\dim(B_{\mathfrak m_1}) = 1$ and $\dim(B_{\mathfrak m_2}) = 2$.
\end{lemma}

\begin{proof}
The construction of $B$ is given above. We just point out that
$B$ is universally catenary by
Algebra, Lemma \ref{algebra-lemma-localization-catenary} and
Morphisms, Lemma \ref{morphisms-lemma-ubiquity-uc}.
\end{proof}

\end{document}
